\documentclass[11pt,a4paper]{article}
\usepackage[utf8]{inputenc}
\usepackage[T1]{fontenc}
\usepackage{amsmath,amssymb,bm,array,geometry,color,fancyhdr,hyperref}
\geometry{left=19mm,right=19mm,top=18mm,bottom=19mm}
\definecolor{ink}{RGB}{20,52,72}
\definecolor{accent}{RGB}{22,115,139}
\hypersetup{colorlinks=true,urlcolor=accent,linkcolor=ink}
\setlength{\parindent}{0pt}\setlength{\parskip}{3pt}
\setlength{\abovedisplayskip}{7pt}\setlength{\belowdisplayskip}{7pt}\setlength{\jot}{4pt}
\pagestyle{fancy}\fancyhf{}\renewcommand{\headrulewidth}{0pt}
\fancyfoot[L]{\small\color{ink}pysnspd / modelo efectivo final / 12 de octubre de 2026}
\fancyfoot[R]{\small\thepage}
\newcommand{\block}[1]{\par\vspace{5pt}{\large\bfseries\color{ink}#1}\par\vspace{3pt}}
\newcommand{\labelrow}[2]{\(\displaystyle #1\)& #2\\}
\newcommand{\sourcecode}[1]{\par{\footnotesize\color{ink}\textbf{Implementación:} \texttt{#1}}\par}
\newcommand{\origin}[1]{\par{\footnotesize\color{ink}\textbf{Campo / procedencia:} #1}\par}
\newcommand{\dd}{\,\mathrm d}
\newcommand{\Eu}{E_u}
\newcommand{\Tb}{T_b}
\newcommand{\Rea}{\operatorname{Re}}
\newcommand{\Ima}{\operatorname{Im}}
\newcommand{\FD}{f_{\rm FD}}
\newcommand{\BE}{n_{\rm BE}}
\newcommand{\Uscale}{u_i^*}
\newcommand{\Rcut}{P_{\rm res}}
\newcommand{\Fgap}{\mathcal F_i}
\newcommand{\meV}{\,\mathrm{meV}}
\begin{document}
{\LARGE\bfseries\color{ink}Ecuaciones principales en 2D}\par
{\large Estado acoplado neutral: \(\Delta,\ p_\xi,\ n_\Omega\)}\par
{\small Modelo efectivo final, etapa 5 completada. Variables físicas en SI; coordenadas espectrales normalizadas explícitas.}
\renewcommand{\theequation}{1.\arabic{equation}}
\block{Condensado: movilidad de Kramer--Watts-Tobin}
\begin{gather}
 d_i=\Delta_i/\Eu,\qquad \Eu=k_BT_c,\qquad
 E_{ik}=\Eu e_{ik},\quad e_{ik}=\sqrt{\bar\xi_k^2+|d_i|^2},\\
 D_td_i=\dot d_i+\frac{2ie}{\hbar}\phi_i d_i,\qquad
 z_i=\begin{cases}d_i/|d_i|,&|d_i|>0,\\1,&d_i=0,\end{cases}\\
 \Fgap=\left.\frac{\partial U_{e+\nabla}}{\partial d_i^*}\right|_{p,\alpha},
 \qquad
 f_i=\frac{z_i^*\Fgap}{\Uscale},\quad \Uscale=N_0\Eu^2V_i,\\
 \boxed{D_td_i=-\frac{z_i}{A_i\tau_K}
 \left[\frac{\Rea f_i}{R_i}+iR_i\Ima f_i\right]},\\
 A_i=\sqrt{\frac{1+\vartheta_i}{2}},\qquad
 \vartheta_i=\frac{\max(T_{E,i},T_b)}{T_c},\qquad
 \tau_K=\frac{\pi\hbar}{8k_BT_c},\\
 \tau_{{\rm sc},i}^{-1}
 =\frac{\vartheta_i}{\tau_{ee}(T_c)}
 +\frac{\vartheta_i^3}{\tau_{ep}(T_c)},\qquad
 R_i=\sqrt{1+\frac{4\Eu^2\tau_{{\rm sc},i}^2|d_i|^2}{\hbar^2}},\\
 \tau_{ee}(T_c)>0,\quad\tau_{ep}(T_c)>0,\qquad
 \bigl(\tau_{ee}(T_c),\tau_{ep}(T_c)\bigr)_{\rm nominal}
 =(5,\ 24.7)\,\mathrm{ps}.
\end{gather}
{\small Calibradores libres: los nominales no fijan el rango ni sustituyen el ajuste experimental.}
\vspace{5pt}
\par\noindent\begin{tabular}{@{}p{.30\linewidth}p{.65\linewidth}@{}}
\hline
\labelrow{i,t,\mathbf r_i}{nodo 2D; tiempo [s]; posición [m]}
\labelrow{\Delta_i,d_i}{parámetro de orden [J]; coordenada adimensional}
\labelrow{E_{ik},\bar\xi_k}{energía de excitación [J]; energía normal nativa dividida por $E_u$}
\labelrow{\phi_i,e,\hbar}{potencial [V]; carga elemental positiva [C]; acción [J\,s]}
\labelrow{U_{e+\nabla},\Fgap}{energía compartida [J]; derivada de Wirtinger [J]}
\labelrow{N_0,V_i,\Uscale}{DOS normal monoespín [J$^{-1}$m$^{-3}$]; volumen [m$^3$]; escala [J]}
\labelrow{T_{E,i},T_b,T_c}{temperatura equivalente electrónica, baño y crítica [K]}
\labelrow{A_i,R_i,f_i,z_i}{factor térmico, factor KWT, fuerza y dirección adimensionales}
\labelrow{\tau_{ee}(T_c),\tau_{ep}(T_c)}{calibradores electrón--electrón y electrón--fonón [s]}
\labelrow{\tau_{\rm sc},\tau_K}{tiempo inelástico combinado; escala cinética KWT [s]}
\hline
\end{tabular}\par
\vfill
\origin{Dinámica disipativa del condensado, teoría KWT; nomenclatura y calibradores de \href{https://thesis.caltech.edu/13748/}{Allmaras (2020)}, ec.\ 3.24, pp.\ 94--98. Fuerza BCS variacional y extrapolación mesoscópica adoptadas aquí.}
\sourcecode{effective\_coupled.py; effective\_bcs\_functional.py}

\newpage
\block{Cuasipartículas y fonones: cinética en estados nativos}
\begin{gather}
 \boxed{\begin{aligned}\dot p_{ik}={}&
 \frac{1}{4w_k}\sum_a\nu^{e}_{ka}r_{ia}
 -\frac{1}{4w_kV_i}\sum_{\ell,a}B_{i\ell}
 \beta^{(i)}_{ka}J_{\ell a}\\
& +\frac{\FD(E_{ik},T_{E,i})-p_{ik}}{\tau_{{\rm BGK},i}}
 +\mathcal H_{ik}\end{aligned}},\\
 \boxed{\dot n_{iq}=
 \frac{1}{c_q}\sum_a\nu^{ph}_{qa}r_{ia}
 -\frac{n_{iq}-\BE(\Omega_q,T_b)}{\tau_{{\rm esc},i}}},\\
 \tau_{{\rm BGK},i}=
 \tau_{ee}(T_c)\frac{T_c}{\max(T_{E,i},T_b)},\qquad
 c_q=\frac{g_{ph}(\Omega_q)W_q^{ph}}{N_0E_u},\\
 0\leq p_{ik}\leq1,\qquad n_{iq}\geq0,\qquad
 \tau_{{\rm esc},i}^{\rm nominal}=15\,\mathrm{ps},\\
 E_{ik}=E_u\sqrt{\bar\xi_k^2+|d_i|^2},\qquad
 \Omega_q=E_u\omega_q,\qquad
 \dot E_{ik}=\frac{E_u\Rea(d_i^*\dot d_i)}{e_{ik}},\\
 \sum_k e_{ik}\nu^e_{ka}
 +\sum_q\omega_q\nu^{ph}_{qa}=0,\qquad
 \sum_q\nu^{ph}_{qa}=1,\\
 B_{i\ell}=\begin{cases}+1,&i=\text{origen de }\ell,\\
 -1,&i=\text{destino de }\ell,\\0,&\text{otro nodo}.\end{cases}
\end{gather}
{\small Fonones locales con escape a interfaz; sin transporte fonónico lateral en este cierre. La energía móvil del gap se contabiliza mediante la energía común, sin advección artificial de etiquetas \(\xi\).}
\vspace{5pt}
\par\noindent\begin{tabular}{@{}p{.30\linewidth}p{.65\linewidth}@{}}
\hline
\labelrow{p_{ik},n_{iq}}{ocupación cuasipartícula y ocupación fonónica adimensionales}
\labelrow{k,q,a,\ell}{estado nativo electrónico; modo fonónico; evento; arista orientada}
\labelrow{w_k,W_q^{ph}}{peso de cuadratura de $\bar\xi$; peso físico fonónico [J]}
\labelrow{c_q,\nu^e,\nu^{ph}}{capacidad nativa; estequiometrías electrón y fonón}
\labelrow{r_{ia}}{extensión neta por segundo del evento local [s$^{-1}$]}
\labelrow{J_{\ell a}}{flujo nativo compartido de transporte [m$^3$/s]}
\labelrow{\beta^{(i)}_{ka},B_{i\ell}}{pesos baricéntricos en energía e incidencia geométrica}
\labelrow{\FD,\BE,T_E}{Fermi--Dirac, Bose--Einstein; funcional de energía equivalente}
\labelrow{\mathcal H_{ik}}{fuente de ocupación por calor físico [s$^{-1}$]}
\labelrow{\tau_{\rm BGK},\tau_{\rm esc}}{relajación conservativa electrónica y escape fonónico [s]}
\labelrow{g_{ph}(\Omega)}{densidad de estados fonónica [J$^{-1}$m$^{-3}$]}
\hline
\end{tabular}\par
\vfill
\origin{Estadística cuántica BCS, ecuaciones cinéticas electrón--fonón, conservación por eventos y volúmenes finitos. \href{https://arxiv.org/pdf/1611.06060}{Vodolazov (2017)}, ecs.\ 1--4 y 13--14; discretización nativa conservativa propia. \(15\) ps es hipótesis de interfaz, no calibración experimental.}
\sourcecode{effective\_bcs\_kinetics.py; effective\_bcs\_kinetics\_numba.py}
\end{document}
